% FILE: 06_contribution_to_thesis.tex
% Project: experiments-visualizer-nextjs
% Role: web-based-pi-visualization-surface

\section{Contribution to the Dissertation}
\label{sec:experiments-visualizer-nextjs-contribution}

\subsection{Contribution to Prediction vs.~Coordination}
\label{sec:experiments-visualizer-nextjs-contribution-prediction}

The dissertation's central axis distinguishes prediction --- the generation of plausible
future states from a model trained on past data --- from coordination --- the alignment of
observed event behaviour against a declared process model in a way that can be audited and
receipted. The visualizer-nextjs project contributes to the coordination side of this axis
by serving as the human-auditable endpoint of the coordination chain.

The CSS design system makes the coordination/prediction distinction visible at the UI layer.
The alignment block panel (\texttt{.alignment-block.sync/.log/.model}) renders the output
of A{*} conformance checking, which is a coordination operation: it compares what happened
(the event log) against what was declared (the Petri net model) and records the moves needed
to reconcile them. The EWMA drift panel (\texttt{.ewma-canvas}, \texttt{.alarm-banner})
renders the output of a statistical process-control operation, which monitors whether the
process is drifting away from its coordinated baseline. Neither of these panels has
anything to do with prediction in the machine-learning sense; both are coordination
evidence surfaces.

AUTHOR THESIS: The existence of a browser-native coordination evidence surface is a
necessary condition for the dissertation's claim that process intelligence can be made
auditable outside specialist tooling. The visualizer-nextjs project is the proposed
implementation of that surface, and its PARTIAL status documents the remaining distance
between specification and proof.

\subsection{Contribution to Process Evidence}
\label{sec:experiments-visualizer-nextjs-contribution-evidence}

Process evidence, in the sense used by the research program, is evidence that a lawful
process occurred --- not merely that a system returned a result. The visualizer-nextjs
project contributes to the process evidence thesis by demonstrating that the type algebra
of process evidence can be encoded in a CSS class vocabulary and used as a rendering
specification for a web dashboard.

Specifically, the \texttt{globals.css} file is a formal enumeration of the observable
states of process evidence across four evidence classes: Petri net marking states, A{*}
alignment move types, EWMA alarm states, and cryptographic receipt verification states.
By encoding these states as CSS classes with no overlap and no ambiguity, the design
system asserts that the visualization surface is type-safe with respect to the process
evidence it displays --- it cannot render an alignment move that is neither sync, log,
nor model; it cannot render a ledger block that is neither healthy nor tampered.

This contribution is genuine even in the project's current PARTIAL state. The design
system is a manufactured artefact that is part of the research record. Its class vocabulary
is cited in this thesis as evidence of what the full visualization surface would need
to encode, and it constitutes a partial specification for the implementation work that
remains.

\subsection{Contribution to Manufacturing Architecture}
\label{sec:experiments-visualizer-nextjs-contribution-manufacturing}

The visualizer-nextjs project contributes to the CodeManufactory manufacturing architecture
by demonstrating that a web-based observation surface can be designed as a manufacturing
target. The \texttt{AGENTS.md} file's instruction to read Next.js 16 documentation from
\texttt{node\_modules/} before writing code reveals that the project was designed with
agentic authorship in mind --- the authoring agent is expected to manufacture React
components against the CSS specification, not hand-author them interactively.

This positions the visualizer's \texttt{globals.css} as a \emph{manufacturing specification}:
a formal description of what the manufacturing pipeline must produce. In the CodeManufactory
framing, the CSS file is the seed; the React component implementations are the bred artefacts;
the test suite is the proof gate; and the receipt is the manufacturing record. The project
is, at PARTIAL status, a case study in a manufacturing pipeline that has produced a valid
seed but has not yet advanced through the subsequent stages.

The project also demonstrates that the manufacturing architecture extends to the web UI layer:
it is not confined to Rust binary artefacts or Python analysis scripts. AUTHOR THESIS: a
complete CodeManufactory proof should include at least one web-surface artefact that passes
all proof gates, and the visualizer-nextjs project is the designated candidate for that role.

\subsection{Contribution to Capital-Flow Infrastructure}
\label{sec:experiments-visualizer-nextjs-contribution-capital}

The dissertation's capital-flow and settlement layer requires an observation surface through
which auditors, counterparties, and governance participants can inspect the process evidence
that underlies capital-allocation decisions. The visualizer-nextjs project is the proposed
web surface for this function.

The cryptographic ledger explorer --- \texttt{.ledger-block-card.healthy/.tampered},
\texttt{.ledger-chain-arrow.broken}, \texttt{.receipt-json-body},
\texttt{.verification-console} --- is the subsystem most directly relevant to capital-flow
infrastructure. In a settlement context, a receipt ledger block represents a unit of
committed process evidence: a claim that a specific process event occurred, was admitted to
the model, was conformance-scored, and was receipted by a tamper-evident chain. The
tamper-detection UI states (\texttt{.tampered}, \texttt{.broken}) make it visible to an
auditor when the evidentiary chain supporting a settlement claim has been compromised.

FUTURE WORK: The full capital-flow use case would require the ledger explorer to display
receipts from multiple counterparty process chains, not just a single local chain, and
to support cross-chain receipt verification. This is beyond the current scope of the
project but is implied by the dashboard grid's design (\texttt{.dashboard-grid} with
a 320-pixel sidebar suitable for a case/counterparty navigator).
